\section{USAMO 1976}
        \begin{enumerate}
        	\item (\textit{USAMO 1976}) 
                    \begin{figure}[H]
                        \centering
                        \includegraphics[width=0.5\textwidth]{Problem/USAMO_I_1976_1_Problem_Figure_1.png}
                    \end{figure}
                    


 
                \begin{itemize}
                    \item (a) Suppose that each square of a $4\times 7$ chessboard, as shown above, is colored either black or white. Prove that with \textit{any} such coloring, the board must contain a rectangle (formed by the horizontal and vertical lines of the board such as the one outlined in the figure) whose four distinct unit corner squares are all of the same color.
\item (b) Exhibit a black-white coloring of a $4\times 6$ board in which the four corner squares of every rectangle, as described above, are not all of the same color.
                \end{itemize}


                



\textbf{Solution 1}

There are many ways to prove the first part, we will show one.

 Consider the first row. It contains at least four cells of the same color, say white. Pick any four such cells. Let's now consider these four columns only. If any of the remaining three rows contains two white cells in some of these columns, we are done. Otherwise, we have the following situation: three rows times four columns, and in each row there is at most one white cell. We can now pick any two rows and then find two columns such that all four cells where they intersect are black, and we are done. (Note that this proof in fact shows that already in a $3\times 7$ board there must be such a rectangle.)

 For the second part, consider the board shown below. Its six columns contain the six different permutations of two white and two black squares. From this observation it is immediately obvious that no two columns agree on two cells of the same color.

 
                    \begin{figure}[H]
                        \centering
                        \includegraphics[width=0.5\textwidth]{Solution/USAMO_I_1976_1_Solution_Figure_1.png}
                    \end{figure}
                    

 


 



\textbf{Solution 2}

We will prove the first part by the pigeonhole principle. To form a rectangle, two squares have in one column have to be colored the same color as the corresponding squares in another column. In the example given, the two second and third squares of the fourth and seventh columns are all white. When coloring this rectangle, every column can be entirely one color, have 3 squares be of one color, or have 2 squares be of one color. Coloring each column with two white squares and two black squares will minimize the number of squares, since a smaller fraction of the column is one color (if more than 2 squares in the same column are the same color, then there is more than one coloring in another column that can create a rectangle). There are ${4}\choose{2}$, or $6$ different ways a column can be colored with two white and two black squares. We need the same coloring in two different columns to create a rectangle, and this is guaranteed with 7 columns by the pigeonhole principle.

 For the second part, we can draw a rectangle with each column having a different coloring with 2 black and 2 white squares. One such rectangle is shown below:

 
                    \begin{figure}[H]
                        \centering
                        \includegraphics[width=0.5\textwidth]{Solution/USAMO_I_1976_1_Solution_Figure_2.png}
                    \end{figure}
                    

 
<i>Alternate solutions are always welcome. If you have a different, elegant solution to this problem, please \href{https://artofproblemsolving.com/wiki/index.php?title=1976\_USAMO\_Problems/Problem\_1&action=edit}{add it} to this page.</i>

 


	\item (\textit{USAMO 1976}) If $A$ and $B$ are fixed points on a given circle and $XY$ is a variable diameter of the same circle, determine the locus of the point of intersection of lines $AX$ and $BY$. You may assume that $AB$ is not a diameter.


 
                    \begin{figure}[H]
                        \centering
                        \includegraphics[width=0.5\textwidth]{Problem/USAMO_I_1976_2_Problem_Figure_1.png}
                    \end{figure}
                    


 



\textbf{Solution}

WLOG, assume that the circle is the unit circle centered at the origin. Then the points $A$ and $B$ have coordinates $(-a,b)$ and $(a,b)$ respectively and $X$ and $Y$ have coordinates $(r,s)$ and $(-r,-s)$. Note that these coordinates satisfy $a^2 + b^2 = 1$ and $r^2 + s^2 = 1$ since these points are on a unit circle. Now we can find equations for the lines:
 \begin{align*} AX \longrightarrow y &= \frac{(s-b)x+rb+sa}{r+a}\\ BY \longrightarrow y &= \frac{(s+b)x+rb-sa}{r+a}. \end{align*}
Solving these simultaneous equations gives coordinates for $P$ in terms of $a, b, r,$ and $s$:$P = \left(\frac{as}{b},\frac{1 - ar}{b}\right)$. These coordinates can be parametrized in Cartesian variables as follows:
 \begin{align*} x &= \frac{as}{b}\\ y &= \frac{1 - ar}{b}. \end{align*}
Now solve for $r$ and $s$ to get $r = \frac{1-by}{a}$ and $s = \frac{bx}{a}$ . Then since $r^2 + s^2 = 1, \left(\frac{bx}{a}\right)^2 + \left(\frac{1-by}{a}\right)^2 = 1$ which reduces to $x^2 + (y-1/b)^2 = \frac{a^2}{b^2}.$ This equation defines a circle and is the locus of all intersection points $P$. In order to define this locus more generally, find the slope of this circle function using implicit differentiation:
 \begin{align*} 2x + 2(y-1/b)y' &= 0\\ (y-1/b)y' &= -x\\ y' &= \frac{-x}{y-1/b}. \end{align*}
Now note that at points $A$ and $B$, this slope expression reduces to $y' = \frac{-b}{a}$ and $y' = \frac{b}{a}$ respectively, values which are identical to the slopes of lines $AO$ and $BO$. Thus we conclude that the complete locus of intersection points is the circle tangent to lines $AO$ and $BO$ at points $A$ and $B$ respectively.

 



\textbf{Solution 2}

Notice that $m\angle AYP=\frac{1}{2}m\overarc{AB}$ (Inscried angle theorem) and that $m\angle XAY=90$ since $XY$ is a diameter, and thus subtends an arc of $180$. This will hold for all $X$ and all $Y$, and so by AA similarity, the angle $APY$ will be constant for all P, thus implying that the points A, B, and all P will be concyclic.If we assume that the center of this circle is $H$, we know that $m\angle AHB=2m\angle APB=2(90-\frac{1}{2}\overarc{AB})=180-\overarc{AB}$. We can assume that $\overline{AH}$ and $\overline{BH}$ intersect the original circle at points $M$, and $N$ respectively. This will give us that $m\angle AHB=\frac{1}{2}|m\overarc{MN}-m\overarc{AB}|$ (Since $H$ lies on the perpendicular bisector of $AB$, we know that $M$, $N$ will be on the same side of $AB$.) Now we also know that $m\overarc{MN}\leq m\overarc{AB}$ or $m\overarc{MN}\leq 360-m\overarc{AB}$. The only case where $m\overarc{MN}$ satisfies the measure of $\angle AHB$, is when $M=A$ and $N=B$, implying that $AH$ and $BH$ are tangents, and so $H$ is the intersection of the tangents from $A$ and $B$ to the original circle. 

 


 <i>Alternate solutions are always welcome. If you have a different, elegant solution to this problem, please \href{https://artofproblemsolving.com/wiki/index.php?title=1976\_USAMO\_Problems/Problem\_2&action=edit}{add it} to this page.</i>

 


	\item (\textit{USAMO 1976}) Determine all integral solutions of $a^2+b^2+c^2=a^2b^2$.


 



\textbf{Solution}

Either $a^2=0$ or $a^2>0$. If $a^2=0$, then $b^2=c^2=0$. Symmetry applies for $b$ as well. If $a^2,b^2\neq 0$, then $c^2\neq 0$. Now we look at $a^2\bmod{4}$:

 $a^2\equiv 0\bmod{4}$: Since a square is either 1 or 0 mod 4, then all the other squares are 0 mod 4. Let $a=2a_1$, $b=2b_1$, and $c=2c_1$. Thus $a_1^2+b_1^2+c_1^2=4a_1^2b_1^2$. Since the LHS is divisible by four, all the variables are divisible by 4, and we must do this over and over again, and from infinite descent, there are no non-zero solutions when $a^2\equiv 0\bmod{4}$.

 $a^2\equiv 1\bmod{4}$: Since $b^2\neq 0\bmod{4}$, $b^2\equiv 1\bmod{4}$, and $2+c^2\equiv 1\bmod{4}$. But for this to be true, $c^2\equiv 3\bmod{4}$, which is an impossibility. Thus there are no non-zero solutions when $a^2\equiv 1\bmod{4}$.

 Thus the only solution is the solution above: $(a,b,c)=0$.

 


 
<i>Alternate solutions are always welcome. If you have a different, elegant solution to this problem, please \href{https://artofproblemsolving.com/wiki/index.php?title=1976\_USAMO\_Problems/Problem\_3&action=edit}{add it} to this page.</i>

 


 


	\item (\textit{USAMO 1976}) If the sum of the lengths of the six edges of a trirectangular tetrahedron $PABC$ (i.e., $\angle APB=\angle BPC=\angle CPA=90^o$) is $S$, determine its maximum volume.


 



\textbf{Solution}

Let the side lengths of $AP$, $BP$, and $CP$ be $a$, $b$, and $c$, respectively. Therefore $S=a+b+c+\sqrt{a^2+b^2}+\sqrt{b^2+c^2}+\sqrt{c^2+a^2}$. Let the volume of the tetrahedron be $V$. Therefore $V=\frac{abc}{6}$.

 Note that $(a-b)^2\geq 0$ implies $\frac{a^2-2ab+b^2}{2}\geq 0$, which means $\frac{a^2+b^2}{2}\geq ab$, which implies $a^2+b^2\geq ab+\frac{a^2+b^2}{2}$, which means $a^2+b^2\geq \frac{(a+b)^2}{2}$, which implies $\sqrt{a^2+b^2}\geq \frac{1}{\sqrt{2}} \cdot (a+b)$. Equality holds only when $a=b$. Therefore

 $S\geq a+b+c+\frac{1}{\sqrt{2}} \cdot (a+b)+\frac{1}{\sqrt{2}} \cdot (c+b)+\frac{1}{\sqrt{2}} \cdot (a+c)$

 $=(a+b+c)(1+\sqrt{2})$.

 $\frac{a+b+c}{3}\geq \sqrt[3]{abc}$ is true from AM-GM, with equality only when $a=b=c$. So $S\geq (a+b+c)(1+\sqrt{2})\geq 3(1+\sqrt{2})\sqrt[3]{abc}=3(1+\sqrt{2})\sqrt[3]{6V}$. This means that $\frac{S}{3(1+\sqrt{2})}=\frac{S(\sqrt{2}-1)}{3}\geq \sqrt[3]{6V}$, or $6V\leq \frac{S^3(\sqrt{2}-1)^3}{27}$, or $V\leq \frac{S^3(\sqrt{2}-1)^3}{162}$, with equality only when $a=b=c$. Therefore the maximum volume is $\frac{S^3(\sqrt{2}-1)^3}{162}$.

 



\textbf{Solution 2}

Note that by AM-GM 
 \[S = a + b + c + \sqrt{a^2 + b^2} + \sqrt{b^2 + c^2} + \sqrt{a^2 + c^2} \ge 3\sqrt[3]{abc} + \sqrt{2ab} + \sqrt{2bc} + \sqrt{2ac} \ge 3\sqrt[3]{abc} + 3\sqrt{2}\sqrt[3]{abc},\]so 
 \[\sqrt[3]{abc} \le \frac{S}{3 + 3\sqrt{2}} = \frac{S}{3} \cdot (\sqrt{2} - 1).\] Proceed as before.

 


 
<i>Alternate solutions are always welcome. If you have a different, elegant solution to this problem, please \href{https://artofproblemsolving.com/wiki/index.php?title=1976\_USAMO\_Problems/Problem\_4&action=edit}{add it} to this page.</i>

 


	\item (\textit{USAMO 1976}) If $P(x)$, $Q(x)$, $R(x)$, and $S(x)$ are all \href{https://artofproblemsolving.com/wiki/index.php?title=Polynomial}{polynomials} such that

 \[P(x^5) + xQ(x^5) + x^2 R(x^5) = (x^4 + x^3 + x^2 + x +1) S(x),\]
prove that $x-1$ is a factor of $P(x)$.


 



\textbf{Solution 1}

In general we will show that if $m$ is an integer less than $n$ and $P_0, \dotsc, P_{m-1}$ and $S$ are polynomials satisfying
 \[P_0(x^n) + x P_1(x^n) + \dotsb + x^{m-1} P_{m-1}(x^n) = \sum_{k=0}^{n-1} x^k S(x),\]then $P_k(1) = 0$, for all integers $0 \le k \le m-1$.  For the problem, we may set $n=5$, $m=3$, and then note that since $P(1)= 0$, $(x-1)$ is a factor of $P(x)$.

 Indeed, let $\omega_1, \dotsc, \omega_{n-1}$ be the $n$th \href{https://artofproblemsolving.com/wiki/index.php?title=Roots\_of\_unity}{roots of unity} other than 1.  Then for all integers $1 \le i \le n-1$,
 \[\sum_{k=0}^{m-1} \omega_i^k P_k(\omega_i^n) = \sum_{k=0}^{m-1} \omega_i^k P_k(1) = \sum_{k=0}^{n-1} \omega_i^k S(\omega_i) = \frac{\omega_i^{n-1} - 1}{\omega_i - 1} S(\omega_i) = 0 ,\]for all integers $1 \le i \le n$.  This means that the $(m-1)$th degree polynomial
 \[\sum_{k=0}^{m-1} x^k P_k(1)\]has $n-1 > m-1$ distinct roots.  Therefore all its coefficients must be zero, so $P_k = 0$ for all integers $0 \le k \le m-1$, as desired.  $\blacksquare$

 



\textbf{Solution 2}

Let $\zeta, \xi, \rho$ be three distinct primitive fifth \href{https://artofproblemsolving.com/wiki/index.php?title=Roots\_of\_unity}{roots of unity}.Setting $x = \zeta, \xi$, we have
 \begin{align*} P(1) + \zeta Q(1) + \zeta^2 R(1) &= \frac{\zeta^5-1}{\zeta-1} S(\zeta) = 0, \\ P(1) + \xi Q(1) + \xi^2 R(1) &= \frac{\xi^5-1}{\xi-1} S(\xi) = 0 . \end{align*}
These equations imply that
 \[\zeta Q(1) + \zeta^2 R(1) = \xi Q(1) + \xi^2 R(1),\]or
 \[Q(1) = - ( \zeta + \xi)R(1).\]But by symmetry,
 \[Q(1) = - (\zeta + \rho)R(1) .\]Since $\xi \neq \rho$, it follows that $Q(1) = R(1) = 0$.  Then, as noted above,
 \[P(1) = P(1) + \zeta Q(1) + \zeta^2 R(1) = 0,\]so $(x-1)$ is a factor of $P(x)$, as desired.  $\blacksquare$

 



\textbf{Solution 3}

Let $z, z^2, z^3$ be three of the 5th roots of unity not equal to one that satisfy $1 + z + z^2 + z^3 + z^4 = 0$ as a result. Plugging them into the equation gives the linear system of equations in $(A(1), B(1), C(1))$:

 
 \[A(1) + zB(1) + z^2C(1) = 0\]
 \[A(1) + z^2B(1) + z^4C(1) = 0\]
 \[A(1) + z^3B(1) + z^6C(1) = 0\]

 Direct observation gives $(A(1), B(1), C(1)) = (0,0,0)$ as a solution, and there are no others because the system of equations is linear. Hence, $A(1) = 0$, and so $(x-1)$ is a factor of $A(x)$, as desired. $\blacksquare$

 \textit{Note: We can generalize this approach to prove the claim in Solution 1 in an equally fast, concise, and readily understandable fashion.}

 


\end{enumerate}